\section{コートアライメント}
\label{sec:alignment}

本節は本稿の中心である。
目標は、再構成シーンの正規化座標と規制コートのメートル座標の間に、
\emph{監査可能で、失敗時に停止する}変換を一つ定めることである。
この変換が決まれば、以降の視点生成と教師生成は解析的な投影だけになる。
逆にここが誤れば、どれだけ高品質なレンダリングも誤った幾何教師を大量生産する。
したがって設計の要は精度の最大化ではなく、
証拠の分離、尺度の一意化、トポロジーの検査、
そして不合格時の明示的な停止である。

\begin{figure}[t]
  \centering
  \resizebox{\linewidth}{!}{% メートル系アライメントの流れ。すべて本稿のために作成した TikZ 図である。
\begin{tikzpicture}[
  font=\sffamily\scriptsize,
  frame/.style={draw=black!28, rounded corners=2.5pt, fill=none,
    inner sep=3mm, line width=0.45pt},
  title/.style={font=\sffamily\bfseries\footnotesize, text=black!85},
  tinylabel/.style={font=\sffamily\fontsize{6.3}{7.2}\selectfont,
    text=black!70, align=center},
  fitcam/.style={regular polygon, regular polygon sides=3, minimum size=3.2mm,
    inner sep=0pt, draw=gsblue, fill=gsblue!20, line width=0.5pt},
  holdcam/.style={regular polygon, regular polygon sides=3, minimum size=3.2mm,
    inner sep=0pt, draw=poseorange, fill=poseorange!22, line width=0.5pt},
  arr/.style={-{Stealth[length=1.7mm,width=1.2mm]}, line width=0.65pt},
  dot/.style={circle, inner sep=0pt, minimum size=1.2mm},
  gate/.style={rounded corners=1.5pt, inner sep=2pt, draw=lossred,
    fill=lossred!8, font=\sffamily\fontsize{6.2}{7.2}\selectfont, align=center},
]

  % ---------- (a) カメラ証拠を fit / holdout に先に分ける ----------
  \begin{scope}[local bounding box=pa]
    \node[fitcam, shape border rotate=-90] (f1) at (0.05,2.05) {};
    \node[fitcam, shape border rotate=-90] (f2) at (1.05,2.25) {};
    \node[fitcam, shape border rotate=-90] (f3) at (2.05,2.05) {};
    \node[holdcam, shape border rotate=-90] (h1) at (1.05,0.05) {};
    \draw[fill=softgray!70, draw=black!30, line width=0.4pt]
      (-0.35,0.60) rectangle (2.45,1.55);
    \node[tinylabel] at (1.05,1.08) {コート面の画像領域};
    \draw[arr, draw=gsblue] (f1) -- (0.40,1.55);
    \draw[arr, draw=gsblue] (f2) -- (1.05,1.55);
    \draw[arr, draw=gsblue] (f3) -- (1.70,1.55);
    \draw[arr, draw=poseorange, dashed] (h1) -- (1.05,0.60);
    \node[tinylabel, text=gsblue, anchor=south] at (1.05,2.50)
      {fit カメラ（変換の推定）};
    \node[tinylabel, text=poseorange, anchor=north] at (1.05,-0.25)
      {holdout カメラ（凍結検証）};
  \end{scope}

  % ---------- (b) 地面推定と光線交差 ----------
  \begin{scope}[shift={(4.6,0)}, local bounding box=pb]
    \coordinate (bl) at (0,-0.30);
    \coordinate (br) at (3.3,-0.30);
    \coordinate (tl) at (0.55,1.05);
    \coordinate (tr) at (2.75,1.05);
    \fill[softgray!80] (bl) -- (br) -- (tr) -- (tl) -- cycle;
    \draw[black!35, line width=0.4pt] (bl) -- (br) -- (tr) -- (tl) -- cycle;
    \draw[black!18] (0.50,-0.30) -- (0.95,1.05)
      (1.15,-0.30) -- (1.45,1.05)
      (1.85,-0.30) -- (1.90,1.05)
      (2.55,-0.30) -- (2.35,1.05);
    \node[fitcam, shape border rotate=-90] (g1) at (0.15,2.20) {};
    \coordinate (e1) at (1.05,0.30);
    \coordinate (e2) at (1.95,-0.20);
    \draw[arr, draw=gsblue] (g1) -- (e1);
    \draw[arr, draw=gsblue, dashed] (g1) -- (e2);
    \node[dot, fill=gsblue] at (e1) {};
    \node[dot, fill=gsblue] at (e2) {};
    \node[tinylabel, text=gsblue, anchor=south west] at (0.45,2.35)
      {線画素の光線と\\推定地上面の交差};
    \node[tinylabel, anchor=north, text width=42mm] at (1.65,-0.55)
      {採用画素数と光線--平面角が\\閾値未満なら、そのカメラを外す};
  \end{scope}

  % ---------- (c) 共通尺度の 2 面当てはめ ----------
  \begin{scope}[shift={(9.7,0)}, local bounding box=pc]
    \draw[courtgreen, line width=0.6pt, fill=courtgreen!7]
      (0.20,0.85) rectangle (1.75,2.05);
    \draw[courtgreen, line width=0.35pt] (0.975,0.85) -- (0.975,2.05);
    \draw[black!28, line width=0.35pt] (0.20,1.45) -- (1.75,1.45);
    \draw[courtgreen, line width=0.6pt, fill=courtgreen!7]
      (2.50,0.55) rectangle (4.05,1.75);
    \draw[courtgreen, line width=0.35pt] (3.275,0.55) -- (3.275,1.75);
    \draw[black!28, line width=0.35pt] (2.50,1.15) -- (4.05,1.15);
    \draw[<->, courtgreen, line width=0.5pt] (0.35,2.20) -- (3.90,2.20);
    \node[tinylabel, text=courtgreen, anchor=south] at (2.125,2.22)
      {共通尺度 \(s\)（メートル換算）};
    \node[gate, anchor=north] at (2.125,0.40)
      {中心間距離 \(\geq 10.97\) m、\\フットプリント重なり \(\leq 10^{-9}\)};
    \node[tinylabel, anchor=north, text width=44mm] at (2.125,-0.55)
      {トポロジーと尺度が食い違えば\\シーン全体を棄却する};
  \end{scope}

  % ---------- パネル枠と見出し ----------
  \node[frame, fit=(pa)] (fa) {};
  \node[frame, fit=(pb)] (fb) {};
  \node[frame, fit=(pc)] (fc) {};
  \node[title, anchor=south] at ($(fa.north)+(0,1.3mm)$) {(a) 証拠の分割};
  \node[title, anchor=south] at ($(fb.north)+(0,1.3mm)$) {(b) 地上面への持ち上げ};
  \node[title, anchor=south] at ($(fc.north)+(0,1.3mm)$) {(c) 共通尺度の 2 面当てはめ};

  \draw[arr] (fa.east) -- (fb.west);
  \draw[arr] (fb.east) -- (fc.west);
\end{tikzpicture}
}
  \caption{\textbf{メートル系アライメントの流れ。}
  カメラは当てはめの前に fit と holdout へ固定分割する。
  fit カメラの線証拠を推定地上面へ持ち上げ、
  共通尺度の規制テンプレートを二面分当てはめ、
  最後に剛体変換をメートル系シーン座標で求める。
  凍結した変換は holdout カメラと全テンプレートのトポロジーを説明しなければならない。
  カメラ端の正準化には二つの固有回転 \(I\) と \(R_z(\pi)\) だけを使う。
  }
  \label{fig:alignment}
\end{figure}

\subsection{証拠の分割と事前検査}
\label{sec:evidence_split}

最初に、どのカメラの証拠を当てはめに使い、
どのカメラを検証に残すかを決める。
この分割は当てはめの前に固定し、結果を見てから入れ替えない。
既定では再構成カメラを復元フレーム番号とカメラ識別子の順に整列し、
時間方向のギャップを二分して得られる入れ子順序から 48 台を選ぶ。
その 48 台を、fit 8 台と holdout 4 台からなる単位で決定的に分割する
（48 台では fit 32 台、holdout 16 台）。
この分割は乱数に依存せず、
変換を求めた後に holdout を当てはめへ戻す経路は実装しない。

事前検査は当てはめの前に走り、
スキーマ、カメラ識別子の一意性、配列形状と型、
有限性、内部パラメータの妥当性、
回転行列の直交性と行列式、
シーン--カメラ変換とその逆変換の整合性を確認する。
ここで落ちたシーンは以降へ進まない。
検出線モデルの読み込みや GPU 確保より前に落とすことで、
設定ミスが長時間の計算の後に発覚することを避ける。

\subsection{地上面の推定と線証拠の持ち上げ}
\label{sec:ground_plane}

線は画像上では二次元の証拠にすぎない。
メートル尺度を得るには、まず線がどの平面上に描かれているかを決める必要がある。
疎な再構成点から候補平面
\(\pi_g=\{\mathbf{x}:\mathbf{n}^{\top}\mathbf{x}+b=0\}\)
を RANSAC で推定し、支持点を三度再当てはめして精緻化する。
採用は支持点数の下限、法線と上方向の余弦、
正のカメラ側の割合、支持点の範囲、
カメラ高さの分布といった複数の検査をすべて満たす場合に限る。

検出した線の画素 \(\tilde{\mathbf{u}}\) に対し、
それを観測した復元カメラはシーン座標の光線
\(\mathbf{r}(\lambda)=\mathbf{o}+\lambda\mathbf{d}\) を定める。
平面との交点は
\begin{equation}
  \lambda^*=-\frac{\mathbf{n}^{\top}\mathbf{o}+b}{\mathbf{n}^{\top}\mathbf{d}}
  \label{eq:ray_plane}
\end{equation}
であり、光線--平面のなす角の余弦、
カメラ前方であること、シーンの範囲内であることの検査を通った点だけを残す。
線の証拠が一つも残らないカメラは証拠集合から明示的に外す。
残った証拠が下限に満たなければ、代替の平面仮定を置かずエラーにする。

持ち上げた点は地上面上の二次元基底で表現し、
メートル尺度未定の証拠集合 \(\mathcal{E}\) とする。
尺度が未定であることが、次の段階の理由である。

\subsection{共通尺度の規制テンプレート当てはめ}
\label{sec:template_fit}

規制コートは寸法が既知である。
そこで正準コートの各線分をサンプルし、
テンプレート \(\mathcal{Q}\) とする。
証拠集合 \(\mathcal{E}\) への距離を
\(d(\mathbf{a},\mathcal{E})=\min_{\mathbf{e}\in\mathcal{E}}\|\mathbf{a}-\mathbf{e}\|_2\)
と書くと、実装した提案スコアは
\begin{equation}
  J(\mathbf{c},\theta,s)=\frac{1}{|\mathcal{Q}|}
  \sum_{\mathbf{q}\in\mathcal{Q}}
  \exp\!\left[-\frac{1}{2}
    \left(\frac{d(s\mathbf{R}(\theta)\mathbf{q}+\mathbf{c},\mathcal{E})}
                    {s\sigma}\right)^{2}\right],
  \qquad
  (\mathbf{c}^*,\theta^*,s^*)=\arg\max J
  \label{eq:template_fit}
\end{equation}
である。ここで \(s\) はメートルあたりのシーン単位（正規化座標の換算）、
スコアの幅は \(\sigma=0.5\) m である。
探索は有界の差分進化法で行い、
タイル化した中心、方位帯、尺度を覆う。

\paragraph{なぜ二面同時に当てはめるか}
一つのコートだけを当てはめると、尺度と位置が縮退し、
見えている一本のベースラインに合うだけの解が残る。
実会場の再構成には複数のコートが写っていることが多いため、
本手法は二面を同時に要求する。
一面目の当てはめ後にその証拠を抑制し、
二面目の仮説を別に探索する。
二つの仮説は、
線トポロジーの整合、フットプリントの非重複（重なり \(\leq 10^{-9}\)）、
中心間距離の下限、元尺度の一致を満たさなければならない。
その後、共通の尺度を固定して中心を再当てはめする。
この制約が、剛体変換を求める前にメートル尺度を一意化する。
\cref{tab:alignment_gates} に採用した閾値を示す。

\begin{table}[t]
  \centering
  \caption{\textbf{採用したアライメント判定の閾値。}
  これらは実行可能な受理条件であり、
  達成された精度ではない。実際の権威は型付き設定ファイル側にあり、
  本表はその要約である。
  完全な一覧は付録~\ref{app:thresholds} に置く。
  }
  \label{tab:alignment_gates}
  \small
  \begin{tabularx}{\linewidth}{@{}p{0.32\linewidth}X@{}}
    \toprule
    判定 & 閾値 \\
    \midrule
    証拠の所有権 & 決定的に選ぶ 48 カメラ。
      fit 8 : holdout 4 の単位で固定分割（48 カメラでは 32 / 16） \\
    地上面 & 支持点 \(\geq 1000\)、
      normal と上方向の余弦 \(\geq 0.97\)、
      正のカメラ側割合 \(=1\) \\
    投影証拠 & 有効なカメラあたり \(\geq 20\) 点 \\
    コート仮説 & ちょうど 2 面。中心間距離 \(\geq 10.97\) m、
      フットプリント重なり \(\leq 10^{-9}\) \\
    共通尺度 & 元尺度の相対不一致 \(\leq 0.07290401\) \\
    対応構築 & 距離 \(\leq 0.25\) m、カメラあたり 3--200 対応 \\
    fit 受理 & カメラ \(\geq 6\)、対応 \(\geq 100\)、
      0.30 m 以内の内点率 \(\geq 90\%\)、RMS と q95 \(\leq 0.30\) m \\
    holdout 受理 & カメラ \(\geq 3\)、対応 \(\geq 80\)、
      0.30 m 以内の内点率 \(\geq 90\%\)、RMS と q95 \(\leq 0.30\) m \\
    全テンプレート & 0.30 m 以内の内点率 \(\geq 60\%\)、
      q95 \(\leq 1.50\) m、全意味セグメントの内点率 \(\geq 50\%\) \\
    変換と投影 & 逆変換の絶対許容 \(10^{-6}\)、
      投影の絶対許容 \(10^{-5}\) px \\
    \bottomrule
  \end{tabularx}
\end{table}


受理された共通尺度を \(s_*\) とすると、メートルアダプタは
\begin{equation}
  \mathbf{x}_{\Scene_n}=s_*\mathbf{x}_{\Scene_m},\qquad
  \mathbf{x}_{\Scene_m}=s_*^{-1}\mathbf{x}_{\Scene_n}
  \label{eq:metric_adapter}
\end{equation}
である。一様尺度だけを置き、剛体変換の中へ尺度を混ぜない。

\subsection{剛体変換とフレームの確定}
\label{sec:kabsch}

正準コートの点 \(\mathbf{x}_j\) と持ち上げたメートル系シーン点
\(\mathbf{y}_j\) の対応が得られたら、最終変換は Kabsch 解
\begin{equation}
  (\mathbf{R}_i,\mathbf{t}_i)=
  \arg\min_{\mathbf{R}\in\SOthree,\mathbf{t}\in\R^3}
    \sum_j \left\|\mathbf{R}\mathbf{x}_j+\mathbf{t}-\mathbf{y}_j\right\|_2^2,
  \quad
  \T{\Scene_m}{\Court_i}=
  \begin{bmatrix}\mathbf{R}_i&\mathbf{t}_i\\\mathbf{0}^{\top}&1\end{bmatrix}
  \label{eq:kabsch}
\end{equation}
である。鏡映が選ばれないよう \(\det\mathbf{R}_i=1\) を強制する。
\cref{eq:kabsch} に尺度が現れないのは、
尺度が \cref{eq:metric_adapter} で既に確定しているためである。

対応は距離 \(\leq 0.25\) m の範囲で構築し、
寄与カメラあたり 3--200 対応の範囲で採用する。
対応点は最近傍だけでなく、
対象コートの物理的な線分上で対応を取る。
これにより、ある一本のベースラインだけに合う解が
全体の姿勢として通ってしまうことを避ける。

\subsection{凍結検証と停止条件}
\label{sec:holdout}

fit カメラの対応から \cref{eq:kabsch} を推定した後、
holdout カメラはこの変換を\emph{再当てはめせずに}評価する。
検証は対応点の最近傍残差だけでなく、
コートテンプレート全体をサンプルして行う。
これは、見えている一本のベースライン上の良い残差が、
大域的に誤ったコートを隠すことを防ぐためである。

受理条件は次の四層からなる。

\begin{enumerate}[leftmargin=1.6em]
  \item \textbf{fit 受理}：カメラ数、対応数、
  \(0.30\) m 以内の内点率、RMS と q95 がすべて閾値内であること。
  \item \textbf{holdout 受理}：同じ指標を凍結変換で満たすこと。
  \item \textbf{全テンプレート}：内点率 \(\geq 60\%\)、
  q95 \(\leq 1.50\) m、
  \emph{すべての}意味セグメントが内点率 \(\geq 50\%\) であること。
  \item \textbf{変換と投影}：相互変換の整合と投影の整合。
\end{enumerate}

第三の条件が重要である。
平均の内点率だけを見ると、
コートの一部が完全にずれていても通ってしまう。
意味セグメントごとの被覆を要求することで、
遠い側のコート、サービスライン、サイドラインのいずれかが
系統的にずれた仮説を棄却する。

すべての条件を満たした変換のみが
\cref{sec:system} の \(\mathcal{L}\) として公開される。
一つでも欠ければシーンは棄却され、
「とりあえず一面だけの変換」や
「holdout を使って再当てはめした変換」を返す経路はない。
この停止は理論上の保証ではなく、
誤った幾何教師が数千枚単位で生成されることを防ぐ実務上の要件である。
過去の B00 の校正試行は、
この停止が実際に発火する例である（当時の判定条件による）。
集約した fit 残差が閾値を満たしても、
fit のカメラ集合を 8 グループずつ三通りに選び直すと
コート中心が最大 \(9.08\) m、方位が \(89.94^\circ\) ずれ、
安定性の判定で棄却された。
その後の局所再当てはめは fit の各 gate を通ったが、
凍結した変換による holdout は距離重み付き q95 が \(1.240\) m、
グループ別のテンプレート被覆が最小 \(0.223\) と下限を下回り、
holdout の判定で棄却された。
いずれの試行でも機械受理のアライメント契約は公開されていない。
これは失敗ではなく、設計どおりの挙動である。

\subsection{カメラ端の正準化}
\label{sec:canonicalization}

コートは線対称に近いため、
カメラがコートのどちら側にあるかで
「近い側」「遠い側」の意味が入れ替わる。
この曖昧性を放置すると、
物理点の番号と姿勢教師が矛盾する。

本研究はコートフレームでカメラ中心の \(Y\) 座標を見て、
二つの固有回転だけを使う。
\begin{equation}
  \T{\Canon}{\Cam}=\mathbf{S}\T{\Court_i}{\Cam},
  \qquad \mathbf{S}\in\{\mathbf{I},\mathbf{R}_z(\pi)\}
  \label{eq:canonical_pose}
\end{equation}
\(Y<0\) なら \(I\)、\(Y>0\) なら \(R_z(\pi)\) を選び、
キーポイントの番号置換と回転を同時に適用する。
|Y| が \(10^{-6}\) m 以下の中央面近傍は、
近い側と遠い側が幾何的に定義できないため、
姿勢で無理に決めずにその sample を棄却する。
左右の割り当ては常にコート局所の \(X\) に従う。

この正準化は\emph{画像から絶対方位を識別した}ことを意味しない。
カメラがコートのどちら側にあるかという物理的な情報を使い、
表現の重複を取り除いているだけである。
